\chapter{Comparacion con otras alternativas}\label{alternativas}

\section{Competidores considerados}

Para el analisis competitivo se han considerado cinco soluciones de referencia en el
mercado espanol de comparacion de compra:

\begin{itemize}
	\item Soysuper
	\item Findit
	\item OCU Market
	\item PreciRadar
	\item RadarPrice
\end{itemize}

Adicionalmente, se han valorado como competencia indirecta las apps nativas de cadenas
y como sustituto no tecnologico la comparacion manual con folletos.

\section{Matriz funcional resumida}

\begin{table}[h]
\centering
\caption{Comparativa de capacidades clave}
\label{tab:comparativa-alternativas}
\begin{tabular}{|p{3.4cm}|p{1.9cm}|p{1.8cm}|p{1.8cm}|p{1.9cm}|}
\hline
\textbf{Capacidad} & \textbf{BargAIn} & \textbf{Soysuper} & \textbf{OCU Market} & \textbf{Resto} \\
\hline
Comparacion multi-tienda de cesta & Completa & Completa & Parcial & Parcial \\
\hline
Optimizacion de ruta (precio+dist+tiempo) & Completa & No & No & No \\
\hline
Integracion de PYMEs/comercio local & Completa & No & No & No \\
\hline
OCR de lista/ticket & Completa & No & No & No \\
\hline
Asistente conversacional IA & Completa & No & No & No \\
\hline
Historico avanzado de precios & Parcial & Limitado & Limitado & Variable \\
\hline
\end{tabular}
\end{table}

\section{Brechas estrategicas detectadas}

El analisis identifica cuatro brechas de mercado que justifican la propuesta:

\begin{enumerate}
	\item \textbf{Gap precio-ruta:} los comparadores muestran precios, pero no optimizan
		desplazamiento y tiempo.
	\item \textbf{Comercio local invisible:} no existe cobertura estable de PYMEs en
		herramientas generalistas.
	\item \textbf{Friccion de entrada:} la lista de compra sigue siendo manual en la mayoria
		de soluciones.
	\item \textbf{Interaccion limitada:} predominan filtros y busquedas rigidas frente a
		consultas en lenguaje natural.
\end{enumerate}

\section{Posicionamiento de BargAIn}

BargAIn se posiciona como un \emph{optimizador inteligente de compra} y no solo como
comparador. Su propuesta de valor combina ahorro economico, eficiencia logistica y
experiencia de usuario asistida por IA.

\section{Riesgos competitivos}

Los riesgos principales identificados son:

\begin{itemize}
	\item reaccion de competidores consolidados con mayor base de usuarios,
	\item dependencia de fuentes de datos de terceros para scraping,
	\item necesidad de mantener alta calidad de datos para sostener confianza.
\end{itemize}

Como mitigacion, se adopta estrategia de fuentes hibridas de datos, modularidad
arquitectonica y pipeline de validacion continua.
	 



























